% Course Study Guide -- shared preamble (all chapters rely on these macros/environments)
% Set the course code in main.tex BEFORE % Course Study Guide -- shared preamble (all chapters rely on these macros/environments)
% Set the course code in main.tex BEFORE % Course Study Guide -- shared preamble (all chapters rely on these macros/environments)
% Set the course code in main.tex BEFORE % Course Study Guide -- shared preamble (all chapters rely on these macros/environments)
% Set the course code in main.tex BEFORE \input{preamble}:  \newcommand{\CourseCode}{ME521}
\providecommand{\CourseCode}{Course}
\usepackage[T1]{fontenc}
\usepackage[utf8]{inputenc}
\usepackage{lmodern}
\usepackage[a4paper,margin=2.4cm,headheight=14pt]{geometry}
\usepackage{microtype}
\usepackage{amsmath,amssymb,amsthm,mathtools}
\usepackage{bm}
\usepackage{enumitem}
\usepackage{booktabs,array,multirow}
\usepackage{graphicx}
\usepackage{xcolor}
\usepackage{tikz}
\usetikzlibrary{arrows.meta,calc,decorations.pathreplacing,patterns,positioning}
\usepackage{pgfplots}
\pgfplotsset{compat=1.17}
\usepackage[most]{tcolorbox}
\usepackage{fancyhdr}
\usepackage{titlesec}
\usepackage[hidelinks]{hyperref}
\usepackage[nameinlink,capitalise]{cleveref}

\definecolor{primary}{HTML}{1F4E79}
\definecolor{defcol}{HTML}{2E75B6}
\definecolor{thmcol}{HTML}{C55A11}
\definecolor{excol}{HTML}{548235}
\definecolor{notecol}{HTML}{7F6000}
\definecolor{warncol}{HTML}{C00000}
\definecolor{solcol}{HTML}{5B2C6F}

\titleformat{\chapter}[display]{\normalfont\huge\bfseries\color{primary}}{\chaptertitlename\ \thechapter}{12pt}{\Huge}
\titleformat{\section}{\normalfont\Large\bfseries\color{primary}}{\thesection}{1em}{}
\titleformat{\subsection}{\normalfont\large\bfseries\color{primary!85!black}}{\thesubsection}{1em}{}

\pagestyle{fancy}
\fancyhf{}
\fancyhead[L]{\small\nouppercase{\leftmark}}
\fancyhead[R]{\small \CourseCode{} Study Guide}
\fancyfoot[C]{\thepage}

\setlist{itemsep=2pt,topsep=4pt}
\numberwithin{equation}{section}

% ---------- boxed environments ----------
% Boxed theorems: \begin{definition}{Title}{label} -> \cref{def:label}; keyidea/pitfall/instructor take [optional title] (commas OK)
\newtcbtheorem[number within=chapter,crefname={Definition}{Definitions},Crefname={Definition}{Definitions}]{definition}{Definition}{enhanced,breakable,colback=defcol!5,colframe=defcol,fonttitle=\bfseries,separator sign={.},description delimiters={(}{)}}{def}
\newtcbtheorem[use counter from=definition,crefname={Theorem}{Theorems},Crefname={Theorem}{Theorems}]{theorem}{Theorem}{enhanced,breakable,colback=thmcol!5,colframe=thmcol,fonttitle=\bfseries,separator sign={.},description delimiters={(}{)}}{thm}
\newtcbtheorem[use counter from=definition,crefname={Example}{Examples},Crefname={Example}{Examples}]{example}{Example}{enhanced,breakable,colback=excol!4,colframe=excol,fonttitle=\bfseries,separator sign={.},description delimiters={(}{)}}{ex}
% NB: tcbtheorems take TWO mandatory args: \begin{theorem}{Title}{label} -- leave either empty: \begin{theorem}{}{}

\newtcolorbox{keyidea}[1][Key idea]{enhanced,breakable,colback=notecol!6,colframe=notecol,fonttitle=\bfseries,title={#1}}
\newtcolorbox{pitfall}[1][Common pitfall]{enhanced,breakable,colback=warncol!4,colframe=warncol,fonttitle=\bfseries,title={#1}}
\newtcolorbox{instructor}[1][From the instructor's notes]{enhanced,breakable,colback=white,colframe=primary!60,borderline west={3pt}{0pt}{primary},fonttitle=\bfseries\itshape,coltitle=primary,colbacktitle=primary!8,title={#1}}
\newtcolorbox{summarybox}[1][Chapter summary]{enhanced,breakable,colback=primary!4,colframe=primary,fonttitle=\bfseries,title={#1}}

% problem / solution for exercise chapters
\newtcolorbox{problem}[1]{enhanced,breakable,colback=gray!6,colframe=gray!60!black,fonttitle=\bfseries,title=Problem #1}
\newenvironment{solution}{\par\noindent\textbf{\color{solcol}Solution.}\ }{\hfill$\blacksquare$\par\medskip}

\theoremstyle{remark}
\newtheorem*{remark}{Remark}

% ---------- math macros ----------
\newcommand{\R}{\mathbb{R}}
\newcommand{\C}{\mathbb{C}}
\newcommand{\N}{\mathbb{N}}
\newcommand{\dd}{\,\mathrm{d}}
\newcommand{\der}[2]{\frac{\mathrm{d}#1}{\mathrm{d}#2}}
\newcommand{\pder}[2]{\frac{\partial #1}{\partial #2}}
\newcommand{\Lap}{\mathcal{L}}          % Laplace transform
\newcommand{\Op}{\mathscr{L}}           % differential operator
\newcommand{\ip}[2]{\left\langle #1,\, #2\right\rangle}
\newcommand{\norm}[1]{\left\lVert #1\right\rVert}
\newcommand{\abs}[1]{\left\lvert #1\right\rvert}
\newcommand{\Wr}{W}                     % Wronskian
\DeclareMathOperator{\sgn}{sgn}
\usepackage{mathrsfs}
:  \newcommand{\CourseCode}{ME521}
\providecommand{\CourseCode}{Course}
\usepackage[T1]{fontenc}
\usepackage[utf8]{inputenc}
\usepackage{lmodern}
\usepackage[a4paper,margin=2.4cm,headheight=14pt]{geometry}
\usepackage{microtype}
\usepackage{amsmath,amssymb,amsthm,mathtools}
\usepackage{bm}
\usepackage{enumitem}
\usepackage{booktabs,array,multirow}
\usepackage{graphicx}
\usepackage{xcolor}
\usepackage{tikz}
\usetikzlibrary{arrows.meta,calc,decorations.pathreplacing,patterns,positioning}
\usepackage{pgfplots}
\pgfplotsset{compat=1.17}
\usepackage[most]{tcolorbox}
\usepackage{fancyhdr}
\usepackage{titlesec}
\usepackage[hidelinks]{hyperref}
\usepackage[nameinlink,capitalise]{cleveref}

\definecolor{primary}{HTML}{1F4E79}
\definecolor{defcol}{HTML}{2E75B6}
\definecolor{thmcol}{HTML}{C55A11}
\definecolor{excol}{HTML}{548235}
\definecolor{notecol}{HTML}{7F6000}
\definecolor{warncol}{HTML}{C00000}
\definecolor{solcol}{HTML}{5B2C6F}

\titleformat{\chapter}[display]{\normalfont\huge\bfseries\color{primary}}{\chaptertitlename\ \thechapter}{12pt}{\Huge}
\titleformat{\section}{\normalfont\Large\bfseries\color{primary}}{\thesection}{1em}{}
\titleformat{\subsection}{\normalfont\large\bfseries\color{primary!85!black}}{\thesubsection}{1em}{}

\pagestyle{fancy}
\fancyhf{}
\fancyhead[L]{\small\nouppercase{\leftmark}}
\fancyhead[R]{\small \CourseCode{} Study Guide}
\fancyfoot[C]{\thepage}

\setlist{itemsep=2pt,topsep=4pt}
\numberwithin{equation}{section}

% ---------- boxed environments ----------
% Boxed theorems: \begin{definition}{Title}{label} -> \cref{def:label}; keyidea/pitfall/instructor take [optional title] (commas OK)
\newtcbtheorem[number within=chapter,crefname={Definition}{Definitions},Crefname={Definition}{Definitions}]{definition}{Definition}{enhanced,breakable,colback=defcol!5,colframe=defcol,fonttitle=\bfseries,separator sign={.},description delimiters={(}{)}}{def}
\newtcbtheorem[use counter from=definition,crefname={Theorem}{Theorems},Crefname={Theorem}{Theorems}]{theorem}{Theorem}{enhanced,breakable,colback=thmcol!5,colframe=thmcol,fonttitle=\bfseries,separator sign={.},description delimiters={(}{)}}{thm}
\newtcbtheorem[use counter from=definition,crefname={Example}{Examples},Crefname={Example}{Examples}]{example}{Example}{enhanced,breakable,colback=excol!4,colframe=excol,fonttitle=\bfseries,separator sign={.},description delimiters={(}{)}}{ex}
% NB: tcbtheorems take TWO mandatory args: \begin{theorem}{Title}{label} -- leave either empty: \begin{theorem}{}{}

\newtcolorbox{keyidea}[1][Key idea]{enhanced,breakable,colback=notecol!6,colframe=notecol,fonttitle=\bfseries,title={#1}}
\newtcolorbox{pitfall}[1][Common pitfall]{enhanced,breakable,colback=warncol!4,colframe=warncol,fonttitle=\bfseries,title={#1}}
\newtcolorbox{instructor}[1][From the instructor's notes]{enhanced,breakable,colback=white,colframe=primary!60,borderline west={3pt}{0pt}{primary},fonttitle=\bfseries\itshape,coltitle=primary,colbacktitle=primary!8,title={#1}}
\newtcolorbox{summarybox}[1][Chapter summary]{enhanced,breakable,colback=primary!4,colframe=primary,fonttitle=\bfseries,title={#1}}

% problem / solution for exercise chapters
\newtcolorbox{problem}[1]{enhanced,breakable,colback=gray!6,colframe=gray!60!black,fonttitle=\bfseries,title=Problem #1}
\newenvironment{solution}{\par\noindent\textbf{\color{solcol}Solution.}\ }{\hfill$\blacksquare$\par\medskip}

\theoremstyle{remark}
\newtheorem*{remark}{Remark}

% ---------- math macros ----------
\newcommand{\R}{\mathbb{R}}
\newcommand{\C}{\mathbb{C}}
\newcommand{\N}{\mathbb{N}}
\newcommand{\dd}{\,\mathrm{d}}
\newcommand{\der}[2]{\frac{\mathrm{d}#1}{\mathrm{d}#2}}
\newcommand{\pder}[2]{\frac{\partial #1}{\partial #2}}
\newcommand{\Lap}{\mathcal{L}}          % Laplace transform
\newcommand{\Op}{\mathscr{L}}           % differential operator
\newcommand{\ip}[2]{\left\langle #1,\, #2\right\rangle}
\newcommand{\norm}[1]{\left\lVert #1\right\rVert}
\newcommand{\abs}[1]{\left\lvert #1\right\rvert}
\newcommand{\Wr}{W}                     % Wronskian
\DeclareMathOperator{\sgn}{sgn}
\usepackage{mathrsfs}
:  \newcommand{\CourseCode}{ME521}
\providecommand{\CourseCode}{Course}
\usepackage[T1]{fontenc}
\usepackage[utf8]{inputenc}
\usepackage{lmodern}
\usepackage[a4paper,margin=2.4cm,headheight=14pt]{geometry}
\usepackage{microtype}
\usepackage{amsmath,amssymb,amsthm,mathtools}
\usepackage{bm}
\usepackage{enumitem}
\usepackage{booktabs,array,multirow}
\usepackage{graphicx}
\usepackage{xcolor}
\usepackage{tikz}
\usetikzlibrary{arrows.meta,calc,decorations.pathreplacing,patterns,positioning}
\usepackage{pgfplots}
\pgfplotsset{compat=1.17}
\usepackage[most]{tcolorbox}
\usepackage{fancyhdr}
\usepackage{titlesec}
\usepackage[hidelinks]{hyperref}
\usepackage[nameinlink,capitalise]{cleveref}

\definecolor{primary}{HTML}{1F4E79}
\definecolor{defcol}{HTML}{2E75B6}
\definecolor{thmcol}{HTML}{C55A11}
\definecolor{excol}{HTML}{548235}
\definecolor{notecol}{HTML}{7F6000}
\definecolor{warncol}{HTML}{C00000}
\definecolor{solcol}{HTML}{5B2C6F}

\titleformat{\chapter}[display]{\normalfont\huge\bfseries\color{primary}}{\chaptertitlename\ \thechapter}{12pt}{\Huge}
\titleformat{\section}{\normalfont\Large\bfseries\color{primary}}{\thesection}{1em}{}
\titleformat{\subsection}{\normalfont\large\bfseries\color{primary!85!black}}{\thesubsection}{1em}{}

\pagestyle{fancy}
\fancyhf{}
\fancyhead[L]{\small\nouppercase{\leftmark}}
\fancyhead[R]{\small \CourseCode{} Study Guide}
\fancyfoot[C]{\thepage}

\setlist{itemsep=2pt,topsep=4pt}
\numberwithin{equation}{section}

% ---------- boxed environments ----------
% Boxed theorems: \begin{definition}{Title}{label} -> \cref{def:label}; keyidea/pitfall/instructor take [optional title] (commas OK)
\newtcbtheorem[number within=chapter,crefname={Definition}{Definitions},Crefname={Definition}{Definitions}]{definition}{Definition}{enhanced,breakable,colback=defcol!5,colframe=defcol,fonttitle=\bfseries,separator sign={.},description delimiters={(}{)}}{def}
\newtcbtheorem[use counter from=definition,crefname={Theorem}{Theorems},Crefname={Theorem}{Theorems}]{theorem}{Theorem}{enhanced,breakable,colback=thmcol!5,colframe=thmcol,fonttitle=\bfseries,separator sign={.},description delimiters={(}{)}}{thm}
\newtcbtheorem[use counter from=definition,crefname={Example}{Examples},Crefname={Example}{Examples}]{example}{Example}{enhanced,breakable,colback=excol!4,colframe=excol,fonttitle=\bfseries,separator sign={.},description delimiters={(}{)}}{ex}
% NB: tcbtheorems take TWO mandatory args: \begin{theorem}{Title}{label} -- leave either empty: \begin{theorem}{}{}

\newtcolorbox{keyidea}[1][Key idea]{enhanced,breakable,colback=notecol!6,colframe=notecol,fonttitle=\bfseries,title={#1}}
\newtcolorbox{pitfall}[1][Common pitfall]{enhanced,breakable,colback=warncol!4,colframe=warncol,fonttitle=\bfseries,title={#1}}
\newtcolorbox{instructor}[1][From the instructor's notes]{enhanced,breakable,colback=white,colframe=primary!60,borderline west={3pt}{0pt}{primary},fonttitle=\bfseries\itshape,coltitle=primary,colbacktitle=primary!8,title={#1}}
\newtcolorbox{summarybox}[1][Chapter summary]{enhanced,breakable,colback=primary!4,colframe=primary,fonttitle=\bfseries,title={#1}}

% problem / solution for exercise chapters
\newtcolorbox{problem}[1]{enhanced,breakable,colback=gray!6,colframe=gray!60!black,fonttitle=\bfseries,title=Problem #1}
\newenvironment{solution}{\par\noindent\textbf{\color{solcol}Solution.}\ }{\hfill$\blacksquare$\par\medskip}

\theoremstyle{remark}
\newtheorem*{remark}{Remark}

% ---------- math macros ----------
\newcommand{\R}{\mathbb{R}}
\newcommand{\C}{\mathbb{C}}
\newcommand{\N}{\mathbb{N}}
\newcommand{\dd}{\,\mathrm{d}}
\newcommand{\der}[2]{\frac{\mathrm{d}#1}{\mathrm{d}#2}}
\newcommand{\pder}[2]{\frac{\partial #1}{\partial #2}}
\newcommand{\Lap}{\mathcal{L}}          % Laplace transform
\newcommand{\Op}{\mathscr{L}}           % differential operator
\newcommand{\ip}[2]{\left\langle #1,\, #2\right\rangle}
\newcommand{\norm}[1]{\left\lVert #1\right\rVert}
\newcommand{\abs}[1]{\left\lvert #1\right\rvert}
\newcommand{\Wr}{W}                     % Wronskian
\DeclareMathOperator{\sgn}{sgn}
\usepackage{mathrsfs}
:  \newcommand{\CourseCode}{ME521}
\providecommand{\CourseCode}{Course}
\usepackage[T1]{fontenc}
\usepackage[utf8]{inputenc}
\usepackage{lmodern}
\usepackage[a4paper,margin=2.4cm,headheight=14pt]{geometry}
\usepackage{microtype}
\usepackage{amsmath,amssymb,amsthm,mathtools}
\usepackage{bm}
\usepackage{enumitem}
\usepackage{booktabs,array,multirow}
\usepackage{graphicx}
\usepackage{xcolor}
\usepackage{tikz}
\usetikzlibrary{arrows.meta,calc,decorations.pathreplacing,patterns,positioning}
\usepackage{pgfplots}
\pgfplotsset{compat=1.17}
\usepackage[most]{tcolorbox}
\usepackage{fancyhdr}
\usepackage{titlesec}
\usepackage[hidelinks]{hyperref}
\usepackage[nameinlink,capitalise]{cleveref}

\definecolor{primary}{HTML}{1F4E79}
\definecolor{defcol}{HTML}{2E75B6}
\definecolor{thmcol}{HTML}{C55A11}
\definecolor{excol}{HTML}{548235}
\definecolor{notecol}{HTML}{7F6000}
\definecolor{warncol}{HTML}{C00000}
\definecolor{solcol}{HTML}{5B2C6F}

\titleformat{\chapter}[display]{\normalfont\huge\bfseries\color{primary}}{\chaptertitlename\ \thechapter}{12pt}{\Huge}
\titleformat{\section}{\normalfont\Large\bfseries\color{primary}}{\thesection}{1em}{}
\titleformat{\subsection}{\normalfont\large\bfseries\color{primary!85!black}}{\thesubsection}{1em}{}

\pagestyle{fancy}
\fancyhf{}
\fancyhead[L]{\small\nouppercase{\leftmark}}
\fancyhead[R]{\small \CourseCode{} Study Guide}
\fancyfoot[C]{\thepage}

\setlist{itemsep=2pt,topsep=4pt}
\numberwithin{equation}{section}

% ---------- boxed environments ----------
% Boxed theorems: \begin{definition}{Title}{label} -> \cref{def:label}; keyidea/pitfall/instructor take [optional title] (commas OK)
\newtcbtheorem[number within=chapter,crefname={Definition}{Definitions},Crefname={Definition}{Definitions}]{definition}{Definition}{enhanced,breakable,colback=defcol!5,colframe=defcol,fonttitle=\bfseries,separator sign={.},description delimiters={(}{)}}{def}
\newtcbtheorem[use counter from=definition,crefname={Theorem}{Theorems},Crefname={Theorem}{Theorems}]{theorem}{Theorem}{enhanced,breakable,colback=thmcol!5,colframe=thmcol,fonttitle=\bfseries,separator sign={.},description delimiters={(}{)}}{thm}
\newtcbtheorem[use counter from=definition,crefname={Example}{Examples},Crefname={Example}{Examples}]{example}{Example}{enhanced,breakable,colback=excol!4,colframe=excol,fonttitle=\bfseries,separator sign={.},description delimiters={(}{)}}{ex}
% NB: tcbtheorems take TWO mandatory args: \begin{theorem}{Title}{label} -- leave either empty: \begin{theorem}{}{}

\newtcolorbox{keyidea}[1][Key idea]{enhanced,breakable,colback=notecol!6,colframe=notecol,fonttitle=\bfseries,title={#1}}
\newtcolorbox{pitfall}[1][Common pitfall]{enhanced,breakable,colback=warncol!4,colframe=warncol,fonttitle=\bfseries,title={#1}}
\newtcolorbox{instructor}[1][From the instructor's notes]{enhanced,breakable,colback=white,colframe=primary!60,borderline west={3pt}{0pt}{primary},fonttitle=\bfseries\itshape,coltitle=primary,colbacktitle=primary!8,title={#1}}
\newtcolorbox{summarybox}[1][Chapter summary]{enhanced,breakable,colback=primary!4,colframe=primary,fonttitle=\bfseries,title={#1}}

% problem / solution for exercise chapters
\newtcolorbox{problem}[1]{enhanced,breakable,colback=gray!6,colframe=gray!60!black,fonttitle=\bfseries,title=Problem #1}
\newenvironment{solution}{\par\noindent\textbf{\color{solcol}Solution.}\ }{\hfill$\blacksquare$\par\medskip}

\theoremstyle{remark}
\newtheorem*{remark}{Remark}

% ---------- math macros ----------
\newcommand{\R}{\mathbb{R}}
\newcommand{\C}{\mathbb{C}}
\newcommand{\N}{\mathbb{N}}
\newcommand{\dd}{\,\mathrm{d}}
\newcommand{\der}[2]{\frac{\mathrm{d}#1}{\mathrm{d}#2}}
\newcommand{\pder}[2]{\frac{\partial #1}{\partial #2}}
\newcommand{\Lap}{\mathcal{L}}          % Laplace transform
\newcommand{\Op}{\mathscr{L}}           % differential operator
\newcommand{\ip}[2]{\left\langle #1,\, #2\right\rangle}
\newcommand{\norm}[1]{\left\lVert #1\right\rVert}
\newcommand{\abs}[1]{\left\lvert #1\right\rvert}
\newcommand{\Wr}{W}                     % Wronskian
\DeclareMathOperator{\sgn}{sgn}
\usepackage{mathrsfs}
